\subsection{Master diagnostic summary}
\label{subsec:master_diagnostic_summary}

The diagnostic pipeline developed above can be summarized as a three-stage
consistency test. First, the untapered regular-core baseline establishes centre
regularity and broad compatibility with the standard energy-condition filters.
Second, the tapered branch suppresses residual boundary density and pressure at
the matching surface, thereby reducing the matching tension that appears in the
untapered branch. Third, curvature and compactness diagnostics test whether the
effective single-function metric remains finite in the regularized core and
separate subcompact, horizon-near, and static-coordinate caution branches.

The preferred fiducial branch is the tapered positive-pressure core
\[
(A,w,\alpha,x_m,q)=(1.0,0.3,0.0,4.0,2),
\]
which is used only as a representative toy-model point. It should not be
interpreted as a unique astrophysical solution. Within the scanned grid, the
tapered branch satisfies NEC, WEC, SEC, DEC, centre regularity, and boundary
softness in all tested cases, while the strict causal-proxy filter selects
240 out of 288 cases. The curvature scan gives finite Ricci-scalar and
Kretschmann-like proxies in all 96 tested cases.

These results support the use of BEY as an effective singularity-replacement
compact-core ansatz. They do not constitute a full Einstein--QCD solution, a
complete Darmois--Israel matching proof, or a nonlinear stability theorem.
